\BourbakiExerciseGroup

\(^{*1}\) Let U denote the additive uniformity on the real line R, and let \(U'\) denote the inverse image of U under the mapping \(x \rightarrow x^{3}\) of R onto itself. Show that \(U'\) is strictly finer than U, but that the Cauchy filters are the same for both uniformities. \(^{*}\)

\begin{enumerate}
\def\labelenumi{\arabic{enumi})}
\setcounter{enumi}{1}
\tightlist
\item
  \begin{enumerate}
  \def\labelenumii{\alph{enumii})}
  \tightlist
  \item
    Let X be an infinite set and let F be a non-trivial ultrafilter on X. Show that if X carries the uniformity \(\mathcal{U}(\mathfrak{F})\) (§ 1, Exercise 1) then the complement of X in its completion \(\hat{X}\) consists of a single point, and that the topological space \(\hat{X}\) can be identified with the space associated with the ultrafilter F (Chapter I, § 6, no. 5, Example).
  \end{enumerate}
\end{enumerate}

*b) Deduce from a) an example of a uniformity on R which is finer than the additive uniformity and with respect to which R is not complete.*

*3) a) Let \(\mathcal{U}\) denote the additive uniformity on the real line \(\mathbf{R}\) , let \(\mathcal{U}_1\) denote the uniformity induced on \(\mathbf{R}\) by that of the extended line \(\overline{\mathbf{R}}\) , and let \(\mathcal{U}_2\) denote the uniformity induced on \(\mathbf{R}\) by that of the one-point compactification \(\tilde{\mathbf{R}} = \mathbf{P}_1(\mathbf{R})\) of \(\mathbf{R}\) . Show that \(\mathcal{U}\) is strictly finer than \(\mathcal{U}_1\) , that \(\mathcal{U}_1\) is strictly finer than \(\mathcal{U}_2\) , and that the topologies induced by these three uniformities are the same; also that the completions of \(\mathbf{R}\) with respect to these three uniformities are respectively \(\mathbf{R}\) , \(\overline{\mathbf{R}}\) and \(\tilde{R}\) . The identity mapping of R extends by continuity to an injective but not surjective mapping \(R \rightarrow \overline{R}\) and to a surjective but not injective mapping \(\overline{R} \rightarrow \tilde{R}\) .

\begin{enumerate}
\def\labelenumi{\alph{enumi})}
\setcounter{enumi}{1}
\tightlist
\item
  Deduce from a) an example of two Hausdorff uniform spaces X, Y and a uniformly continuous bijection \(u: X \rightarrow Y\) whose extension by continuity \(\overline{u}: \hat{X} \rightarrow \hat{Y}\) is neither injective nor surjective.*
\end{enumerate}

¶ 4) With the notation of Exercise 5 of § 2, suppose that all the uniform spaces \(X_i\) are complete. Show that the space \(X\) is complete with respect to the uniformity there defined (argue by reductio ad absurdum, using § 2, Proposition 5, Corollary 2).

\begin{enumerate}
\def\labelenumi{\arabic{enumi})}
\setcounter{enumi}{4}
\item
  Let X be a Hausdorff uniform space such that X is the intersection of a countable family of open sets of \(\hat{X}\) . Show that if Y is any Hausdorff uniform space which contains X as a uniform subspace, then X is the intersection of a closed subset of Y and a countable family of open sets of Y.
\item
  Let X be a uniform space and let \(i: X \to \hat{X}\) be the canonical mapping of X into its Hausdorff completion. Let \(X_{0}\) be the sum of X and \(\hat{X} - i(X)\) , and let \(j: X_{0} \to \hat{X}\) be the mapping which agrees with i on X and with the identity mapping on \(\hat{X} - i(X)\) . Show that \(X_{0}\) is complete with respect to the uniformity which is the inverse image under j of the uniformity of \(\hat{X}\) , also that if Y is any complete uniform space such that X is a uniform subspace of Y, then the identity mapping \(X \to X\) can be extended uniquely to a continuous mapping \(X_{0} \to Y\) .
\end{enumerate}

¶ 7) Let \(\mathbf{X}\) be a Hausdorff uniform space, \(\mathcal{U}\) its uniformity, \(\mathcal{U}_0\) the uniformity induced on \(\mathfrak{F}(\mathbf{X})\) by the uniformity \(\tilde{\mathcal{U}}\) ( \(\S\) 2, Exercise 6 b), \(\mathcal{U}_{00}\) the uniformity induced on \(\mathfrak{F}(\mathfrak{F}(\mathbf{X}))\) by the uniformity \(\tilde{\mathcal{U}}_0\) . Show that the canonical mapping \(x \to \{\{x\}\}\) of \(\mathbf{X}\) into \(\mathfrak{F}(\mathfrak{F}(\mathbf{X}))\) extends to an isomorphism of \(\hat{\mathbf{X}}\) onto a closed uniform subspace of \(\mathfrak{F}(\mathfrak{F}(\mathbf{X}))\) (carrying the uniformity \(\mathcal{U}_{00}\) ) (use Proposition 9 of no. 4).

